We study decision-focused learning (DFL) for the discrete Berth Allocation Problem at the Port of San Antonio, Chile. A predictor estimates each vessel's berth service time, and these estimates feed a mixed-integer linear program that schedules a week of berth assignments. Because the predicted service time enters the precedence constraints rather than a linear objective cost vector, the Smart Predict-then-Optimize surrogate does not apply, and we differentiate through the solver using black-box and perturbed gradients. Aligning the DFL loss and the regret metric with the solver's unweighted objective yields an exact non-negative-regret guarantee. Our central hypothesis is cascade asymmetry: under-predicting a service time propagates downstream delays, whereas over-predicting only idles a berth, so a decision-focused loss should outperform mean squared error when schedules are contention-sensitive. On controlled synthetic instances DFL reduces schedule regret, and the benefit grows with berth contention, reaching about a third lower mean regret at moderate contention while remaining negligible at low contention. On real 2019--2025 vessel-call data DFL does not reliably improve on predict-then-optimize: the effect lies within run-to-run variance, and the black-box gradient is unstable, making early stopping essential. We contribute the first decision-focused formulation of berth allocation in the predicted-in-constraints regime, a loss-regret-objective alignment with a non-negative-regret guarantee, and an honest characterization of when DFL helps and of the obstacles that remain.
